 \documentclass[12pt]{article}
 % LaTeX file created by GUIDE version 46.2 on 04/30/26 at 09:50
 %%%% Usage: latex tree_90.tex, dvips tree_90.tex, ps2pdf tree_90.tex
 %%%% or: xelatex tree_90.tex
 \usepackage{pstricks}
 %%%% If using pdflatex, replace above line with \usepackage[pdf]{pstricks}
 %%%% and use: pdflatex -shell-escape tree_90.tex
 \usepackage{pstricks,pst-node,pst-tree}
 \usepackage{geometry}
 \usepackage{xcolor}
 \usepackage{lscape}
 \definecolor{wheat}{rgb}{0.96, 0.87, 0.7}
 \definecolor{lightgray}{rgb}{0.9, 0.9, 0.9}
 \definecolor{aqua}{rgb}{0.0, 1.0, 1.0}
 \definecolor{mauve}{rgb}{0.88, 0.69, 1.0}
 \definecolor{skyblue}{RGB}{86,180,233}
 \definecolor{vermillion}{RGB}{213,94,0}
 \definecolor{purple}{RGB}{204,121,167}
 \definecolor{bluishgreen}{RGB}{0,158,115}
 \pagestyle{empty}
 \begin{document}
 %\begin{landscape}
 \begin{center}
\psset{linecolor=black,tnsep=1pt,tndepth=0cm,tnheight=0cm,treesep=1.5cm,levelsep=55pt,radius=10pt,fillstyle=solid}
  \pstree[treemode=D]{\Tcircle{ 1 }~[tnpos=l]{\shortstack[r]{\texttt{\detokenize{X17p_status}}\\=\texttt{\detokenize{Loss}}}}
 }{
    \Tcircle[fillcolor=orange]{ 2 }~{\shortstack[c]{\emph{70}\\11.34}}  
  \pstree[treemode=D]{\Tcircle[linecolor=brown]{ 3 }~[tnpos=l]{\shortstack[r]{\texttt{\detokenize{V1}}\\=\texttt{\detokenize{Untreated}}}}
 }{
  \pstree[treemode=D]{\Tcircle[linecolor=gray]{ 6 }~[tnpos=l]{\shortstack[r]{\texttt{\detokenize{V2}}\\$\leq_*$\texttt{97.20}}}
 }{
    \Tcircle[fillcolor=yellow]{\small 12}~{\shortstack[c]{\emph{217}\\44.12}}  
    \Tcircle[fillcolor=yellow]{\small 13}~{\shortstack[c]{\emph{393}\\29.95}}  
   }
  \pstree[treemode=D]{\Tcircle{ 7 }~[tnpos=l]{\shortstack[r]{\texttt{\detokenize{CLL_EPITYPE}}\\=\texttt{\detokenize{n_CLL}}}}
 }{
    \Tcircle[fillcolor=orange]{\small 14}~{\shortstack[c]{\emph{48}\\ 0.30}}  
    \Tcircle[fillcolor=yellow]{\small 15}~{\shortstack[c]{\emph{280}\\30.66}}  
 }
 }
 }
 \end{center}
GUIDE v.46.2 0.250-SE
 piecewise-constant  0.100-quantile regression tree
for predicting \texttt{\detokenize{surv_5yrs}}.
 % Tree constructed from 1008 observations
 % (excluding observations with non-positive weight or missing values
 % in \texttt{D}, \texttt{T}, \texttt{R} or \texttt{Z} variables)
 % and pruned from a tree with 226 terminal nodes.
 % Maximum number of split levels is 15 and minimum node sample size is 2.
At each split, an observation goes to the left branch 
 if and only if the condition is satisfied.
 Symbol `$\leq_*$' stands for `$\leq$ or missing'.
\texttt{\detokenize{V1}} = \texttt{\detokenize{TREATMENT_STATUS_AT_SAMPLING}}.
\texttt{\detokenize{V2}} = \texttt{\detokenize{IGHV_IDENTITY_PERCENTAGE}}.
 Splits at nodes drawn with \textbf{\textcolor{gray}{gray}} circles are not statistically significant.
 Splits at nodes drawn with \textbf{\textcolor{brown}{brown}} circles are interaction splits.
 Sample size (in \emph{italics}) and 0.100-quantile of \texttt{\detokenize{surv_5yrs}} printed below nodes.
 Terminal nodes with quantiles above and below value of 24.49 at root node are painted \colorbox{yellow}{yellow} and \colorbox{orange}{orange} respectively.
 Second best split variable at root node is \texttt{\detokenize{TREATMENT_STATUS}}.
 %\end{landscape}
 \end{document}
